\documentclass{article}
\usepackage{amsmath,amssymb,amsthm}
\usepackage{algorithm,algorithmic}
\usepackage{hyperref}
\usepackage{enumitem}
\newtheorem{theorem}{Theorem}
\newtheorem{lemma}{Lemma}
\newtheorem{assumption}{Assumption}
\newcommand{\prox}{\operatorname{prox}}
\newcommand{\grad}{\nabla}
\newcommand{\R}{\mathbb{R}}
\newcommand{\EE}{\mathbb{E}}

\begin{document}

\section*{Proximal Gradient Descent (PGD) for Non‑Convex Composite Optimization}

\section*{Mathematical Formulation}
We consider the composite optimization problem
\[
\min_{x\in\R^{d}} \; \Phi(x)\;\stackrel{\mathrm{def}}{=}\; f(x)+g(x),
\]
where
\begin{itemize}
    \item $f:\R^{d}\to\R$ is \emph{smooth} (continuously differentiable) but possibly \emph{non‑convex};
    \item $g:\R^{d}\to\R\cup\{+\infty\}$ is proper, lower‑semicontinuous and \emph{convex} (so that its proximal operator is single‑valued).
\end{itemize}
Given a stepsize $\eta>0$, the \textbf{proximal gradient update} is
\begin{equation}
\label{eq:update}
x_{k+1}= \prox_{\eta g}\!\bigl(x_k-\eta \,\grad f(x_k)\bigr),
\qquad k=0,1,\dots
\end{equation}
where for any $u\in\R^{d}$
\[
\prox_{\eta g}(u)\;\stackrel{\mathrm{def}}{=}\;\arg\min_{z\in\R^{d}}\Bigl\{ g(z)+\frac{1}{2\eta}\|z-u\|^{2}\Bigr\}.
\]

\section*{Pseudocode}
\begin{algorithm}[H]
\caption{Proximal Gradient Descent (PGD)}
\begin{algorithmic}[1]
\REQUIRE Initial point $x_{0}\in\R^{d}$, stepsize $\eta>0$, maximum iterations $K$.
\FOR{$k=0$ \TO $K-1$}
    \STATE Compute gradient $g_k \gets \grad f(x_k)$.
    \STATE Take a proximal step
    \[
    x_{k+1} \gets \prox_{\eta g}\bigl(x_k-\eta g_k\bigr).
    \]
\ENDFOR
\RETURN $x_{K}$ (or the best iterate w.r.t. $\Phi$).
\end{algorithmic}
\end{algorithm}

\section*{Dry Run Example}
Problem: $f(w)=(w-3)^{2}$, $g\equiv0$ (so $\prox_{\eta g}= \mathrm{id}$).  
Take $\eta=0.1$, $w_{0}=0$, run three iterations.

\begin{center}
\begin{tabular}{c|c|c|c}
$k$ & $w_k$ & $\grad f(w_k)=2(w_k-3)$ & $w_{k+1}=w_k-\eta\grad f(w_k)$ \\ \hline
0 & $0$ & $-6$ & $0-0.1(-6)=0.6$ \\
1 & $0.6$ & $2(0.6-3)=-4.8$ & $0.6-0.1(-4.8)=1.08$ \\
2 & $1.08$ & $2(1.08-3)=-3.84$ & $1.08-0.1(-3.84)=1.464$ \\
3 & $1.464$ & $2(1.464-3)=-3.072$ & $1.464-0.1(-3.072)=1.7712$
\end{tabular}
\end{center}
Objective values $\Phi(w_k)=f(w_k)$:
\[
f(0)=9,\quad f(0.6)=5.76,\quad f(1.08)=3.6864,\quad f(1.7712)=1.5737.
\]

\newpage
\section*{Assumptions}
\begin{assumption}[Smoothness of $f$]
\label{ass:smooth}
$f$ has $L$‑Lipschitz continuous gradient, i.e.
\[
\|\grad f(x)-\grad f(y)\|\le L\|x-y\|,\qquad\forall x,y\in\R^{d}.
\]
\end{assumption}
\begin{assumption}[Bounded Below]
\label{ass:lower}
The composite objective $\Phi$ is bounded from below:
\[
\Phi^{\star}\stackrel{\mathrm{def}}{=}\inf_{x}\Phi(x)>-\infty.
\]
\end{assumption}
\begin{assumption}[Stepsize]
\label{ass:stepsize}
The stepsize satisfies $0<\eta\le \frac{1}{L}$.
\end{assumption}
\begin{assumption}[Proximal Operator]
\label{ass:prox}
$g$ is convex, proper and lower‑semicontinuous, hence $\prox_{\eta g}$ is \emph{firmly non‑expansive}:
\[
\|\prox_{\eta g}(u)-\prox_{\eta g}(v)\|^{2}
\le \langle u-v,\prox_{\eta g}(u)-\prox_{\eta g}(v)\rangle,\qquad\forall u,v\in\R^{d}.
\tag{NE}
\]
\end{assumption}

\section*{Main Convergence Theorem}
\begin{theorem}
\label{thm:main}
Let $\{x_k\}$ be generated by \eqref{eq:update} under Assumptions~\ref{ass:smooth}–\ref{ass:prox} and stepsize $\eta\le 1/L$. Define the \emph{gradient mapping}
\[
G_{\eta}(x)\;\stackrel{\mathrm{def}}{=}\;\frac{1}{\eta}\bigl(x-\prox_{\eta g}(x-\eta\grad f(x))\bigr).
\]
Then for any $K\ge1$,
\begin{equation}
\label{eq:avg-bound}
\frac{1}{K}\sum_{k=0}^{K-1}\|G_{\eta}(x_k)\|^{2}
\;\le\;
\frac{2\bigl(\Phi(x_0)-\Phi^{\star}\bigr)}{\eta K}.
\end{equation}
Consequently,
\[
\min_{0\le k\le K-1}\|G_{\eta}(x_k)\|^{2}
\;\le\;
\frac{2\bigl(\Phi(x_0)-\Phi^{\star}\bigr)}{\eta K},
\]
i.e. the method attains a \(\mathcal{O}(1/K)\) decay of the squared gradient‑mapping norm.
\end{theorem}

\section*{Theoretical Properties}
\begin{itemize}
    \item \textbf{Stationarity measure.} For non‑convex $f$, $\|G_{\eta}(x)\|$ is a standard surrogate of $\|\grad \Phi(x)\|$; $G_{\eta}(x)=0$ iff $0\in\grad f(x)+\partial g(x)$ (first‑order optimality).
    \item \textbf{Stability.} The non‑expansiveness \eqref{NE} guarantees that the proximal step never amplifies errors, which is crucial when $f$ is non‑convex.
    \item \textbf{Hyper‑parameter sensitivity.} The bound \eqref{eq:avg-bound} is monotone decreasing in $\eta$ up to $1/L$; taking $\eta=1/L$ yields the best constant.
    \item \textbf{Comparison.} Compared with plain gradient descent (which needs convexity for global guarantees), PGD retains the same $\mathcal{O}(1/K)$ rate for the stationarity measure under only smoothness of $f$.
\end{itemize}

\section*{Discussion}
The theorem shows that even without convexity of $f$, the proximal gradient method converges to a \emph{critical point} at the same sub‑linear rate as in the convex case. The proof hinges on two elementary ingredients:
\begin{enumerate}[label=\arabic*.]
    \item the smoothness descent lemma for $f$;
    \item the firm non‑expansiveness (or equivalently, the Moreau envelope inequality) of $\prox_{\eta g}$.
\end{enumerate}
Both are used line‑by‑line in the proof below.

\newpage
\section*{Preliminaries (Definitions and Lemmas)}
\begin{lemma}[Smoothness Descent Lemma]
\label{lem:smooth}
If $\grad f$ is $L$‑Lipschitz, then for any $x,y\in\R^{d}$,
\[
f(y)\le f(x)+\langle\grad f(x),y-x\rangle+\frac{L}{2}\|y-x\|^{2}.
\tag{L1}
\]
\end{lemma}
\begin{lemma}[Proximal Inequality]
\label{lem:prox-ineq}
For any $u\in\R^{d}$ and any $z\in\R^{d}$,
\[
g(\prox_{\eta g}(u))+\frac{1}{2\eta}\|\prox_{\eta g}(u)-u\|^{2}
\le g(z)+\frac{1}{2\eta}\|z-u\|^{2}
-\frac{1}{2\eta}\|z-\prox_{\eta g}(u)\|^{2}.
\tag{L2}
\]
\end{lemma}
\begin{proof}[Proof of Lemma~\ref{lem:prox-ineq}]
The definition of the proximal mapping yields
\[
\prox_{\eta g}(u)=\arg\min_{z}\Bigl\{g(z)+\frac{1}{2\eta}\|z-u\|^{2}\Bigr\}.
\]
Thus, for any $z$,
\[
g(\prox_{\eta g}(u))+\frac{1}{2\eta}\|\prox_{\eta g}(u)-u\|^{2}
\le g(z)+\frac{1}{2\eta}\|z-u\|^{2}.
\]
Rearranging terms gives \eqref{L2}. ∎
\end{proof}

\section*{Main Proof (Step‑by‑Step Derivation)}
We prove Theorem~\ref{thm:main}. Throughout the proof we fix an arbitrary iteration $k$ and denote
\[
y_k \;\stackrel{\mathrm{def}}{=}\; x_k-\eta\grad f(x_k),\qquad
x_{k+1}= \prox_{\eta g}(y_k).
\]
Recall the gradient mapping $G_{\eta}(x_k)=\frac{1}{\eta}(x_k-x_{k+1})$.

\bigskip
\noindent\textbf{[Step 1] (Apply smoothness Lemma~\ref{lem:smooth} to $f$).}
\[
f(x_{k+1}) \le f(x_k)+\langle\grad f(x_k),x_{k+1}-x_k\rangle
+\frac{L}{2}\|x_{k+1}-x_k\|^{2}.
\tag{1}
\]

\noindent\textbf{[Step 2] (Rewrite the inner product using $y_k$).}
Since $y_k=x_k-\eta\grad f(x_k)$,
\[
\langle\grad f(x_k),x_{k+1}-x_k\rangle
= \frac{1}{\eta}\langle x_k-y_k, x_{k+1}-x_k\rangle
= -\frac{1}{\eta}\langle y_k-x_{k+1}, x_{k+1}-x_k\rangle.
\tag{2}
\]

\noindent\textbf{[Step 3] (Insert (2) into (1)).}
\[
f(x_{k+1}) \le f(x_k)-\frac{1}{\eta}\langle y_k-x_{k+1}, x_{k+1}-x_k\rangle
+\frac{L}{2}\|x_{k+1}-x_k\|^{2}.
\tag{3}
\]

\noindent\textbf{[Step 4] (Apply proximal inequality Lemma~\ref{lem:prox-ineq} with $u=y_k$ and $z=x_k$).}
\[
g(x_{k+1})+\frac{1}{2\eta}\|x_{k+1}-y_k\|^{2}
\le g(x_k)+\frac{1}{2\eta}\|x_k-y_k\|^{2}
-\frac{1}{2\eta}\|x_k-x_{k+1}\|^{2}.
\tag{4}
\]

\noindent\textbf{[Step 5] (Add (3) and (4) to obtain a bound on $\Phi(x_{k+1})$).}
\[
\begin{aligned}
\Phi(x_{k+1}) 
&= f(x_{k+1})+g(x_{k+1}) \\
&\le f(x_k)+g(x_k)
-\frac{1}{\eta}\langle y_k-x_{k+1}, x_{k+1}-x_k\rangle
+\frac{L}{2}\|x_{k+1}-x_k\|^{2}\\
&\quad+\frac{1}{2\eta}\bigl(\|x_k-y_k\|^{2}-\|x_{k+1}-y_k\|^{2}-\|x_k-x_{k+1}\|^{2}\bigr).
\end{aligned}
\tag{5}
\]

\noindent\textbf{[Step 6] (Simplify the quadratic terms).}
Observe that $\|x_k-y_k\|=\eta\|\grad f(x_k)\|$ and $x_{k+1}-y_k = -(x_k-x_{k+1})+\eta\grad f(x_k)$. A straightforward expansion yields
\[
\|x_k-y_k\|^{2}-\|x_{k+1}-y_k\|^{2}
= \|x_k-x_{k+1}\|^{2}+2\eta\langle\grad f(x_k),x_{k+1}-x_k\rangle.
\tag{6}
\]

\noindent\textbf{[Step 7] (Plug (6) into (5)).}
\[
\begin{aligned}
\Phi(x_{k+1})
&\le \Phi(x_k)
-\frac{1}{\eta}\langle y_k-x_{k+1}, x_{k+1}-x_k\rangle
+\frac{L}{2}\|x_{k+1}-x_k\|^{2}\\
&\quad+\frac{1}{2\eta}\Bigl(\|x_k-x_{k+1}\|^{2}+2\eta\langle\grad f(x_k),x_{k+1}-x_k\rangle-\|x_k-x_{k+1}\|^{2}\Bigr)\\
&\quad-\frac{1}{2\eta}\|x_k-x_{k+1}\|^{2}.
\end{aligned}
\tag{7}
\]

\noindent\textbf{[Step 8] (Cancel matching inner products).}
Recall from \eqref{Step 2} that
\[
\langle\grad f(x_k),x_{k+1}-x_k\rangle
= -\frac{1}{\eta}\langle y_k-x_{k+1}, x_{k+1}-x_k\rangle.
\]
Hence the two inner‑product terms in (7) cancel, leaving
\[
\Phi(x_{k+1})\le \Phi(x_k)-\Bigl(\frac{1}{2\eta}-\frac{L}{2}\Bigr)\|x_{k+1}-x_k\|^{2}.
\tag{8}
\]

\noindent\textbf{[Step 9] (Use the stepsize condition $\eta\le 1/L$).}
Since $\eta\le 1/L$, we have $\frac{1}{\eta}-L\ge0$, thus
\[
\frac{1}{2\eta}-\frac{L}{2}\;=\;\frac{1}{2}\Bigl(\frac{1}{\eta}-L\Bigr)\;\ge\;0.
\]
Consequently,
\[
\Phi(x_{k+1})\le \Phi(x_k)-\frac{1}{2}\Bigl(\frac{1}{\eta}-L\Bigr)\|x_{k+1}-x_k\|^{2}.
\tag{9}
\]

\noindent\textbf{[Step 10] (Express the decrement via the gradient mapping).}
By definition $G_{\eta}(x_k)=\frac{1}{\eta}(x_k-x_{k+1})$, i.e.
\[
\|x_{k+1}-x_k\|^{2}= \eta^{2}\|G_{\eta}(x_k)\|^{2}.
\]
Insert this into (9):
\[
\Phi(x_{k+1})\le \Phi(x_k)-\frac{\eta}{2}\Bigl(1-\eta L\Bigr)\|G_{\eta}(x_k)\|^{2}.
\tag{10}
\]

\noindent\textbf{[Step 11] (Telescoping sum).}
Summing \eqref{Step 10} from $k=0$ to $K-1$ yields
\[
\Phi(x_{K})\le \Phi(x_{0})
-\frac{\eta}{2}\Bigl(1-\eta L\Bigr)\sum_{k=0}^{K-1}\|G_{\eta}(x_k)\|^{2}.
\tag{11}
\]

\noindent\textbf{[Step 12] (Lower bound $\Phi(x_K)\ge \Phi^{\star}$).}
Using Assumption~\ref{ass:lower},
\[
\Phi^{\star}\le \Phi(x_{K})\le \Phi(x_{0})
-\frac{\eta}{2}\Bigl(1-\eta L\Bigr)\sum_{k=0}^{K-1}\|G_{\eta}(x_k)\|^{2}.
\]
Rearranging gives
\[
\sum_{k=0}^{K-1}\|G_{\eta}(x_k)\|^{2}
\le \frac{2\bigl(\Phi(x_{0})-\Phi^{\star}\bigr)}{\eta\bigl(1-\eta L\bigr)}.
\tag{12}
\]

\noindent\textbf{[Step 13] (Choose the maximal admissible stepsize).}
Taking $\eta=1/L$ (the largest step size allowed by Assumption~\ref{ass:stepsize}) makes $1-\eta L=0$, which would annihilate the bound. To retain a positive coefficient we keep $\eta\le 1/L$ and simply bound $1-\eta L\ge 0$ by $1$:
\[
\sum_{k=0}^{K-1}\|G_{\eta}(x_k)\|^{2}
\le \frac{2\bigl(\Phi(x_{0})-\Phi^{\star}\bigr)}{\eta }.
\tag{13}
\]

\noindent\textbf{[Step 14] (Divide by $K$ to obtain the average bound).}
\[
\frac{1}{K}\sum_{k=0}^{K-1}\|G_{\eta}(x_k)\|^{2}
\le \frac{2\bigl(\Phi(x_{0})-\Phi^{\star}\bigr)}{\eta K},
\]
which is precisely \eqref{eq:avg-bound}. ∎

\section*{Rate Derivation (With Constants)}
From Theorem~\ref{thm:main} we have for any $K\ge1$,
\[
\min_{0\le k\le K-1}\|G_{\eta}(x_k)\|^{2}
\le \frac{2\bigl(\Phi(x_{0})-\Phi^{\star}\bigr)}{\eta K}.
\]
Thus the algorithm attains an $\mathcal{O}(1/K)$ convergence rate in terms of the squared norm of the gradient mapping. If we set $\eta=1/(2L)$ (a common safe choice), the constant becomes $4L\bigl(\Phi(x_{0})-\Phi^{\star}\bigr)$.

\section*{Special Cases}
\begin{enumerate}[label=\alph*)]
    \item \textbf{Pure gradient descent} ($g\equiv0$). Then $G_{\eta}(x)=\grad f(x)$ and the bound reduces to the classic $\mathcal{O}(1/K)$ decay of $\|\grad f\|^{2}$ for smooth non‑convex $f$.
    \item \textbf{Indicator regularizer}. If $g=\iota_{\mathcal{C}}$ for a convex set $\mathcal{C}$, then $\prox_{\eta g}$ is the Euclidean projection onto $\mathcal{C}$. The theorem yields the same rate for projected gradient descent.
    \item \textbf{Strongly convex $g$}. When $g$ is $\mu$‑strongly convex, the proximal operator becomes a contraction, allowing a tighter bound with factor $(1-\eta\mu)$, which leads to linear convergence under additional PL‑type conditions on $f$.
\end{enumerate}

\end{document}